% GENERATED, edit the French unit or the translation table
\mxnewpage\section*{Solutions EV3-A}\label{EV3-A-CORR}
\textbf{MCQ:} 1 A, 2 B, 3 A, 4 C, 5 C, 6 D, 7 B, 8 D, 9 D, 10 B, 11 A, 12 B. Each correct answer: 0.5 point; total 6.
\paragraph{Checks of the fluency questions.} Order of operations; sum of fractions; subtraction of exponents; multiplying coefficient 1.2; division by 0.8; product 1.25 times 0.8; solving an affine equation; vertical increment divided by horizontal; mean; angle divided by 360; complement to 1; substitution in brackets.
\subsection*{Exercise 1 -- Lottery and dispersion}
\begin{enumerate}
\item X is -3,1,9 with the same probabilities. \textit{Marking: 1 point. Result and basic justification expected; half a point possible for a correct but unfinished approach.}
\item $E(X)=-\dfrac{3}{2}+\dfrac{3}{8}+\dfrac{9}{8}=0$. $E(X^2)=\dfrac{9}{2}+\dfrac{3}{8}+\dfrac{81}{8}=15$; variance 15. \textit{Marking: 2 points. Method or justification: 1; calculation and conclusion: 1.}
\item Net gain 2X, expectation 0; standard deviation $2\sqrt{15}$. \textit{Marking: 2 points. Method or justification: 1; calculation and conclusion: 1.}
\end{enumerate}
\subsection*{Exercise 2 -- Derivative and horizontal tangent}
\begin{enumerate}
\item $f(0)=-3$; $f\prime=(2x-1)e^x$, $f\prime(0)=-1$. \textit{Marking: 1 point. Result and basic justification expected; half a point possible for a correct but unfinished approach.}
\item Decreasing up to $\dfrac{1}{2}$, increasing after; minimum $-2e^{\dfrac{1}{2}}$. \textit{Marking: 2 points. Method or justification: 1; calculation and conclusion: 1.}
\item $y=-x-3$; contact (0;-3), below the horizontal axis. \textit{Marking: 2 points. Method or justification: 1; calculation and conclusion: 1.}
\end{enumerate}
\subsection*{Exercise 3 -- Angle and length}
\begin{enumerate}
\item AB=4, $AC=\sqrt{4+12}=4$. \textit{Marking: 1 point. Result and basic justification expected; half a point possible for a correct but unfinished approach.}
\item Product 8; cosine $\dfrac{8}{16}=\dfrac{1}{2}$, angle $\dfrac{\pi}{3}$. \textit{Marking: 1 point. Result and basic justification expected; half a point possible for a correct but unfinished approach.}
\item $BC=\sqrt{4+12}=4$: equilateral triangle. \textit{Marking: 2 points. Method or justification: 1; calculation and conclusion: 1.}
\end{enumerate}
\paragraph{Calibration.} Target duration: 20--25 min for the fluency questions, 80--85 min for the exercises, 10--15 min for proofreading. These durations are editorial estimates to be measured in class. Accept any correct method; do not penalise twice a numerical error that is propagated consistently. The proposed marking scheme is not a decision of the team to harmonise marking.
